\section{The eight guarantees}
\label{sec:guarantees}

Throughout this section, \Head{} and the configuration are fixed and
$\st_1 \leq \st_2 \leq \cdots$ is a run sequence (\cref{def:seq}). Two
hypotheses are stated once and used by name.

\begin{description}[leftmargin=3.4em,style=nextline,font=\normalfont\bfseries]
  \item[(H1) stable inherited-channel state] every baseline channel read by
    \code{Propagate-Channel: inherit} remains unchanged through the sequence.
  \item[(H2) no suppressed source is released] for every unit $u$ and every
    $p \in \scopef{u}$ whose contribution to $u$ is suppressed at $\st_1$,
    whether by cancellation, by a hold (\cref{def:sources}), or by a correction
    (\cref{rem:corrections}), $p$ is not released at \Head{} in the
    sequence.
\end{description}

The \emph{retry invariant} of CCME 2.0.0 additionally fixes, for every
outstanding unit, the corrected tuple stream, source sets, graph, scope and
depth predicates, target baselines, and resolved source/target channel
admission. The only permitted state changes are baseline tags written for
packages that successfully release. Thus it entails the constancies used
below; (H1) and (H2) name two historically problematic components, not a
weaker substitute for the full invariant.

Both clauses are explicit parts of CCME 2.0.0's retry invariant. They are
named because the 1.0.0 proof audit isolated them in counterexamples
(\cref{ob:origin,ob:g5}); they are not sufficient on their own. The 2.0.0
fresh window resolves the prerelease-discharge defect of \cref{ob:pre}, while
retries still require every other input listed above to remain fixed. Every
theorem below states when it needs that full invariant.

\subsection{Three lemmas}

\begin{lemma}[window antitonicity]
\label{lem:window}
Let $\st' = \Run{\st}{F}$. For every $p \in \Pkgs$:
\[
  \Win{\st'}{p} =
  \begin{cases}
    \emptyset & \text{if $p$ published a non-prerelease version in the run,}\\
    \Win{\st}{p} & \text{otherwise.}
  \end{cases}
\]
In particular $\Win{\st'}{p} \subseteq \Win{\st}{p}$, and along a run sequence
$\Win{\st_j}{p} \subseteq \Win{\st_k}{p}$ for $j \geq k$.
For the undelivered window, every successful release of $p$, stable or
prerelease, gives $\Winf{\st'}{p}=\emptyset$; otherwise it is unchanged.
\end{lemma}

\begin{proof}
By \cref{def:run} the run adds only tags of the form $(p, \vers{\st}{p},
\Head)$. If $p$ is not among them, its tag set is unchanged, hence so are
$\stcommit{\st}{p}$ and its window. Suppose $p$ is among them and
$\vers{\st}{p}$ is not a prerelease. Then $\vers{\st}{p}$ exceeds
$\base{\st}{p}$ in SemVer precedence, and therefore exceeds
$\stbase{\st}{p}$: \spec{13.9} computes it as $\applyb{\stbase{\st}{p}}{b}$
for the graduation and stable cases, and the guards \diag{E185}, \diag{E195},
\diag{E153} and \diag{E156} reject any computed or pinned version that fails to
exceed the baseline. So the new tag is the highest-precedence non-prerelease
tag of $p$, giving $\stcommit{\st'}{p} = \Head$ and
$\Win{\st'}{p} = \Rch{\Head}\setminus\Rch{\Head} = \emptyset$. If
$\vers{\st}{p}$ is a prerelease, it is not a candidate for $\stbase{}$ at all,
so $\stcommit{}$ and the train window are unchanged. Its tag is nevertheless
the new baseline at \Head{}, so $\Winf{}$ is empty. The final claim follows by
induction along the sequence.
\end{proof}

\begin{remark}
The third case of the proof, a prerelease release leaving the window
untouched, is not an oversight in \cref{def:window} but its stated purpose
(\cref{rem:stable-window}). Admission nevertheless uses the empty undelivered window. \Cref{ob:pre} records the 1.0.0 defect that motivated this separation.
\end{remark}

\begin{lemma}[outstanding-input stability]
\label{lem:source}
Assume the retry invariant. For every outstanding unit $u$ and unreleased
target $d$, its corrected tuple, source set, graph and scope/depth predicates,
target baseline, resolved source and target channels, and channel-eligibility
decision are the same in every retry.
\end{lemma}

\begin{proof}
This is the retry invariant specialised to one unit and target. H1 prevents an
inherited baseline channel from drifting; H2 prevents a successful source
release from changing discharge and hence the source set. The remaining
constancies are separate clauses of the invariant and cannot be derived from
H1 and H2 alone.
\end{proof}

\begin{lemma}[window restriction]
\label{lem:channel}
Under the retry invariant, admission of an outstanding unit at an unreleased
target can change only by restriction of the target's undelivered window. A
successful release of the target empties that window.
\end{lemma}

\begin{proof}
All other admission inputs are fixed by the retry invariant. The first assertion
then follows from \cref{def:admission}; the second is \cref{lem:window}.
\end{proof}

\begin{remark}[why H1 is part of the invariant]
\label{rem:origin-drift}
Without H1, the origin channel of an \code{inherit} unit can change when a
source baseline moves from a train to \stable{}. That can change an unreleased
dependent's resolved channel and planned version between retries.
\Cref{ob:origin} preserves the 1.0.0 counterexample; CCME 2.0.0 prevents it by
fixing inherited baseline-channel state in the retry invariant.
\end{remark}

\subsection{G1: termination}

\begin{theorem}[\G{1}, termination]
\label{thm:g1}
Every traversal of \cref{def:reach} halts after at most $|\Pkgs|$ insertions
into \code{seen} and $O(|\Pkgs| + |\Edges|)$ work, regardless of the declared
depth, including $\infty$ (\code{+*}), and regardless of whether the graph
is acyclic. Propagation performs at most $2|\Units|$ such traversals, so it
halts, in $O(|\Units|\cdot(|\Pkgs|+|\Edges|))$.
\end{theorem}

\begin{proof}
A package is inserted into \code{seen} only after a test that it is absent, and
\code{seen} is never cleared, so there are at most $|\Pkgs|$ insertions. Only
inserted packages (and the initial sources, inserted at seeding) ever enter
\code{next} and hence \code{frontier}, so each package is expanded at most
once, and the edges scanned across the whole traversal are therefore at most
$\sum_{p}|\Dep{p}| = |\Edges|$. The loop terminates either by exhausting
$\ell$ iterations or, when $\ell = \infty$, by the empty-frontier break, which
must occur within $|\Pkgs|$ iterations since each non-empty frontier consumes
at least one distinct unvisited package. No step of this argument uses
acyclicity: a cycle merely means some edge leads to a package already in
\code{seen}, and that edge is skipped. The outer loops of \spec{9.2} run over
the finite set \Units{} once per axis.
\end{proof}

\begin{remark}
Acyclicity (\cref{def:graph}) is therefore not what makes propagation
terminate; \code{seen} is. Acyclicity is needed for the publish order
(\cref{prop:order}), not here. This matters because it means an implementation
that loosens the \diag{E200} check does not silently gain a non-terminating
traversal; it gains an unorderable publish phase, which is a different and
louder failure.
\end{remark}

\subsection{G2: completeness}

\begin{theorem}[\G{2}, completeness: no orphans]
\label{thm:g2}
Assume the retry invariant. If $\admit{\st_1}{u}{d}$, then for every $j$, either $d$
was released in some run $i < j$, or $\admit{\st_j}{u}{d}$. Consequently
$\propb{\st_j}{d} \geq \pbump{u}$ in every run until $d$ releases at a commit
whose ancestry contains $\cof{u}$.
\end{theorem}

\begin{proof}
Suppose $d$ has not been released in runs $1..j-1$. Take the conjuncts of
\cref{def:admission} in turn. \Cref{a:live} depends only on the unit.
\Cref{a:scope} and \cref{a:cancel} depend only on \Head{} and the
configuration. \Cref{a:reach} is $d \in
\reachset{\srcp{\st_j}{u}}{\pdepth{u}}$, and by the retry invariant
the source set is unchanged, as are the graph and the depth, so the traversal
is literally the same computation. \Cref{a:window} is $\cof{u} \in
\Winf{\st_j}{d}$; $d$ gained no tag, so by \cref{lem:window} its window is
unchanged. \Cref{a:resolvable} is constant by \cref{lem:channel}. So every
conjunct still holds.

For the final claim: $\propb{\st_j}{d}$ is a join over the admitted units
(\cref{def:propagated}) and $u$ is among them, so the join is at least
$\pbump{u}$. The only way \cref{a:window} can fail is for $\cof{u}$ to leave
$\Winf{\st_j}{d}$, which by \cref{lem:window} requires $d$ to be released at
\Head{}; and $\cof{u} \anc \Head$, so that release carries $\cof{u}$ in its
ancestry.
\end{proof}

The theorem is the formal content of ``there are no orphans''. What it excludes
is the failure \spec{13.7a} describes: a provider publishes, a consumer's
publish fails, and the consumer is thereafter permanently one patch behind a
dependency it declares, with its manifest range never reconciled and no
diagnostic fired, because nothing appears to be pending. Note where the proof
would break under the other reading of \cref{a:window}. If admission tested
$\cof{u} \in \Win{\st_j}{p}$ for the unit's own package $p$, then $p$'s
successful release would empty \emph{that} window, the unit would cease to be a
propagation source, and $d$ would never be released, in this run or any later
one. Nothing else in the model changes; that one substitution is the whole
difference between a resumable release and a silently lossy one.

\subsection{G3: version stability}

\begin{theorem}[\G{3}, version stability]
\label{thm:g3}
Assume the retry invariant. If $d \in \plan{\st_1}$ and $d$ is not released in runs
$1..k-1$, then $d \in \plan{\st_k}$ and $\vers{\st_k}{d} = \vers{\st_1}{d}$.
\end{theorem}

\begin{proof}
$d$ gained no tag, so $\base{\st_k}{d} = \base{\st_1}{d}$ and
$\stbase{\st_k}{d} = \stbase{\st_1}{d}$, and by \cref{lem:window} its window is
unchanged. Then $\dirb{\st_k}{d} = \dirb{\st_1}{d}$, because
\cref{def:direct} reads the window, $\scopef{u}$, and $\srcp{\st}{u}$, the last
constant by the retry invariant. And $\propb{\st_k}{d} = \propb{\st_1}{d}$: the
$\geq$ direction is \cref{thm:g2} applied to each unit admitted in run $1$, and
the $\leq$ direction is \cref{thm:g5}. Hence $\effb{\st_k}{d} = \effb{\st_1}{d}$. Its resolved and baseline channels, changed-provider status, and applicable exact \code{Release-As} cause are also fixed by the retry invariant. Thus whichever disjunct of \cref{def:plan} put $d$ in the initial plan still holds. For a fixed version function, \spec{13.9} reads unchanged stable and current baselines, resolved channel, applicable exact version, and corrected train records. Those inputs are fixed for an unreleased target by the retry invariant, so its result is unchanged. This proves replay stability. The graduation core itself follows the specific train-target rule of \spec{11.5}; \cref{ob:graduation} records the inconsistent \spec{13.9} summary rather than treating that intent as absent.
\end{proof}

\begin{remark}[the prerelease case survives, for a different reason]
\label{rem:g3-pre}
For a target $d$ planned on a prerelease channel, the version is
$\coreof{\cdot}$-plus-counter per \spec{11.4}, and the counter is read from
$\base{\st}{d}$, which is unchanged precisely because $d$ was \emph{not}
released. This preserves the counter, but extending \cref{thm:g3} to
prerelease targets still requires invariant resolved channels and admitted
contributions, including H1's channel-stability component. The
prerelease problem of \cref{ob:pre} is about packages that \emph{were}
released, and is a convergence failure, not a stability failure.
\end{remark}

\Cref{thm:g3} is as much an operational property as a formal one: after a
partial failure, the re-run publishes the same numbers the operator already
reviewed. It is what makes a resumed release a continuation rather than a fresh
decision.

\begin{example}[four theorems in one table]
\label{ex:runtable}
\Cref{tab:runs} continues \cref{ex:orphan} through an adversarial schedule:
run 1 loses \code{cli}, run 2 loses it again, run 3 succeeds. Each column is a
theorem. \code{docs-site} and \code{@acme/theme} never appear, at any distance
$2$, in any run: the unit said \code{\^{}}, and \cref{thm:g5} forbids a later
run from reconsidering that.
\end{example}

\begin{table}[t]
\centering\small
\begin{tabular}{@{}L{1.1cm}L{3.3cm}L{2.6cm}L{4.4cm}L{2.4cm}@{}}
\toprule
Run & State \st{} & Plan $\plan{\st}$ & Published & Empty? \\
\midrule
1 & $\st_0$ (\cref{ex:workspace})
  & \code{core} 1.5.0, \code{cli} 2.0.1, \code{ui} 0.9.2, \code{api} 1.2.1
  & \code{core}, \code{ui}, \code{api}; \code{cli} fails & no \\
2 & $\st_1 = \st_0 \cup 3$ tags
  & \code{cli} 2.0.1
  & nothing; \code{cli} fails again & no \\
3 & $\st_2 = \st_1$
  & \code{cli} 2.0.1
  & \code{cli} & no \\
4 & $\st_3 = \st_2 \cup \{\code{cli@2.0.1}\}$
  & $\emptyset$ & --- & yes \\
\bottomrule
\end{tabular}
\caption{The run sequence of \cref{ex:runtable}. \Cref{thm:g2}: \code{cli} is
in every plan until it publishes. \Cref{thm:g3}: it is at \code{2.0.1} in all
of them. \Cref{thm:g4}: \code{core}, \code{ui} and \code{api} leave after run
1. \Cref{thm:g5}: nothing new ever enters. \Cref{thm:g6}: the plan shrinks to
empty. Run 2 publishes nothing, which \cref{thm:g6} treats as the hard-error
case rather than the partial-failure one.}
\label{tab:runs}
\end{table}

\subsection{G4: no double release}

\begin{theorem}[\G{4}, no double ledger delivery]
\label{thm:g4}
If $d$ is released in run $k$ at a commit whose ancestry contains $\cof{u}$,
then $\admit{\st_j}{u}{d}$ is false for every $j>k$.
\end{theorem}

\begin{proof}
The successful tag is at \Head{} and contains $\cof{u}$ in its ancestry, so
$\cof{u}\notin\Winf{\st_j}{d}$. Clause~\ref{a:window} of admission therefore
fails. The full train window may retain $\cof{u}$ for aggregate prerelease
version arithmetic; it cannot re-admit the contribution.
\end{proof}

\begin{corollary}[exactly-once ledger discharge]
\label{cor:exactly-once}
Under the retry invariant, a contribution admitted to $d$ remains outstanding until $d$ releases (\cref{thm:g2}) and cannot be delivered in the tag ledger again after that release (\cref{thm:g4}). If retries eventually publish $d$, it is therefore discharged exactly once in that ledger.
\end{corollary}

This is the idempotent-consumer property of the saga, obtained without a
deduplication table: the git tag is the record, and the window is the query
against it. This is not an exactly-once guarantee for external side effects.
A publish can succeed before its tag is recorded; recovery across that gap
requires the verified adoption of \spec{19.4} or an idempotent publisher.
The ledger proof alone cannot prevent a repeated external operation.

\subsection{G5: no blast-radius widening}

\begin{theorem}[\G{5}, no blast-radius widening]
\label{thm:g5}
Assume the retry invariant. For every unit $u$ and every $j \geq 1$,
\[
  \{\, d : \admit{\st_j}{u}{d} \,\} \;\subseteq\; \{\, d : \admit{\st_1}{u}{d} \,\}.
\]
\end{theorem}

\begin{proof}
Compare the conjuncts of \cref{def:admission}, as in \cref{thm:g2}, but now in
the direction of restriction. \Cref{a:live,a:scope,a:cancel} are constant.
\Cref{a:reach} is constant by the retry invariant: the traversal has
the same source set, the same depth and the same graph, all fixed at \Head{}.
\Cref{a:resolvable} is constant under the retry invariant; for a target that has been released \cref{a:window} fails by \cref{lem:window}. \Cref{a:window} is antitone by
\cref{lem:window}. A conjunction of constant and antitone predicates is
antitone.
\end{proof}

\begin{remark}[why this is a security property]
\spec{18.1} depends on \cref{thm:g5}. If the admitted set could grow between
runs, an attacker able to induce a publish failure (by exhausting a
registry quota, say) could enlarge the set of packages a commit releases
\emph{after} that commit had been reviewed and merged. The theorem says the
achievable effect of inducing a failure is confined to delaying a subset of
what was already going to be released.
\end{remark}

\subsection{G6: convergence}

\begin{theorem}[\G{6}, convergence]
\label{thm:g6}
Assume the retry invariant, and suppose every run whose plan is nonempty publishes at least one package. Let
$n = |\plan{\st_1}|$. Then $\plan{\st_{n+1}} = \emptyset$.
\end{theorem}

\begin{proof}
First, $\plan{\st_{j+1}} \subseteq \plan{\st_j}$. Take $p \in
\plan{\st_{j+1}}$. If $p$ was released in run $j$, its new baseline at
\Head{} empties $\Winf{\st_{j+1}}{p}$, so no old direct, propagated, channel,
or exact-version cause can re-admit it under the invariant. Thus it was not
released, and by \cref{lem:window} its window is unchanged
and its tags are unchanged. Then $\dirb{\st}{p}$ is unchanged, $\propb{\st}{p}$ can
only have shrunk (\cref{thm:g5} applied to each unit), $\chanres{\st}{p}$ and
$\chanof{\base{\st}{p}}$ are unchanged under the retry invariant,
and the \code{Release-As} disjunct is unchanged; so every disjunct of
\cref{def:plan} that holds at $\st_{j+1}$ held at $\st_j$, and $\held{p}$ is
constant.

Now let $\mu_j = |\plan{\st_j}|$. Run $j$ publishes some $p \in \plan{\st_j}$,
whose undelivered work is absent from $\plan{\st_{j+1}}$, while by the
paragraph above nothing is added. So $\mu_{j+1} < \mu_j$. A strictly
decreasing sequence in $\mathbb{N}$ bounded below by $0$ reaches $0$ within
$\mu_1 = n$ steps.
\end{proof}

The bound is restricted to the invariant domain. Publishing a suppressed
source or changing channel admission may expose a finite follow-up obligation;
\spec{19.6} requires re-planning and renewed review, not an unconditional
global-convergence claim. Held packages remain excluded.

\begin{proposition}[universal tagging is necessary, \spec{13.10a}]
\label{prop:tagall}
Assume the retry invariant. Suppose the protocol were weakened so that a released package $p$ whose publish
target uploads no artefact is versioned and has its manifest written, but is
not tagged. Then if $\effb{\st_1}{p} \neq \bnone$, $p \in \plan{\st_j}$ for
every $j$, and \cref{thm:g6} fails.
\end{proposition}

\begin{proof}
$p$ gains no tag in any run, so by \cref{lem:window} $\Winf{\st_j}{p} =
\Winf{\st_1}{p}$ for all $j$, and $\dirb{\st_j}{p} = \dirb{\st_1}{p} \neq
\bnone$ if the bump was direct (and, by \cref{thm:g2}, $\propb{\st_j}{p} \geq
\propb{\st_1}{p}$ if it was propagated, since $p$ never releases at \Head{} in
the sense \cref{lem:window} requires). Either way the plan's first disjunct
holds at every $\st_j$, so $\mu_j \geq 1$ for all $j$ and no run empties the
plan.
\end{proof}

This is why \spec{13.10a} tags whatever is released, including private,
internal and artefact-less packages, and why the convergence check of
\spec{19.6} (\diag{E199}) admits exactly one exclusion, held packages. A
permanent exception to convergence would not be a caveat to document; it would
disable the one check that detects the orphan failure of \spec{13.7a}, and
disable it precisely for the internal-only packages where nobody is watching a
registry.

\subsection{G7: channel discharge}

\begin{theorem}[\G{7}, channel discharge]
\label{thm:g7}
Suppose $d$ is released in run $k$ on channel $ch$, so that
$\chanof{\base{\st_j}{d}} = ch$ for $j > k$. Then for every unit $u$ and every
$j > k$, $\mathrm{pcf}_{\st_j}(u,d) \neq ch$: the channel $d$ has arrived on is
never proposed for it again.
\end{theorem}

\begin{proof}
Write $cur = \chanof{\base{\st_j}{d}} = ch$; this is read from $d$'s tag, not
from anything computed in the run. Every branch of \cref{def:propchan} ends in
the test ``if $target = cur$, return \NONE''. A result of $ch$ would require
$target = ch = cur$, which that test rejects. The transition branch fails one
step earlier where it applies: having arrived at $t$, $d$ no longer matches an
$f \neq t$, and for $f = *$ with $t = \stable$ the match requires
$cur \neq \stable$, which fails.
\end{proof}

\begin{corollary}[common ledger discharge]
\label{cor:g7}
After any release, $\Winf{}$ excludes every delivered directive. The no-op and
transition checks of \cref{thm:g7} independently prevent re-proposing the
channel already recorded in the baseline. These are complementary checks;
the baseline comparison is not a substitute for fresh-work admission.
\end{corollary}

Both axes therefore discharge through the common undelivered window, with the
channel axis adding its no-op check. Under the retry invariant both survive a
partial failure unchanged. This also marks the boundary of what
\cref{thm:g7} gives. It says the channel $d$ arrived on is not re-proposed; it
does \emph{not} say nothing is proposed. The historical CCME 1.0.0 counterexample in \cref{ob:pre} used its unsplit window to let a second, older directive propose a different channel on every run. CCME 2.0.0's fresh-window gate discharges that directive after release.

\subsection{G8: axis independence}

\begin{theorem}[\G{8}, axis independence]
\label{thm:g8}
The channel assignment $\chanres{\st}{\cdot}$ is a function of \st{}, \Head{}
and the configuration that does not read any quantity computed by phase 3.
Consequently, altering the bump directives $\pbump{u}, \pdepth{u},
\pscope{u}$ of any set of units leaves $\chanres{\st}{\cdot}$ unchanged, and in
particular a \diag{W208} suppression removes no channel assignment and causes
no package to lose a channel change it would otherwise have had.
\end{theorem}

\begin{proof}
By \cref{def:phase1}, phase 1 reads the units' channel directives, $\srcp{\st}{u}$,
the traversal, $\Winf{\st}{\cdot}$, $\canc{\cdot}{\cdot}$, and, through
\cref{def:propchan}, $\base{\st}{\cdot}$. By \cref{def:channels}, phase 2
reads phase 1's output and the same fixed quantities. None of these is a
phase-3 output: $\propb{\st}{\cdot}$, $\effb{\st}{\cdot}$ and the admission
predicate appear nowhere in \cref{def:propchan,def:phase1,def:channels}. So
$\chanres{\st}{\cdot}$ is determined before phase 3 runs and is invariant under
any change confined to phase-3 inputs.
\end{proof}

The dependency between the axes is thus a single directed edge, from channels
to bumps, and \cref{thm:g8} is the statement that there is no edge back. Two
operational consequences follow, both worth stating because they are what a
plan reader needs: a \diag{W208} in a plan means a caret did not reach, not
that a channel directive failed; and a package may legitimately appear in a
plan with a channel change and no bump (\diag{W202}), or with a bump and no
channel change, and neither shape indicates a defect.

\begin{remark}[the converse does not hold]
\label{rem:g8-converse}
CCME 1.0.0 stated \G{8} as ``suppressing a bump under \spec{9.3a} never
suppresses a channel, \emph{and vice versa}''. CCME 2.0.0 states only the
one-way result above. The old proof observed that
the passes share only $\chanres{\st}{\cdot}$, which phase 3 reads and never writes.
That argument establishes exactly the direction proved above. The converse is
false, and intentionally so: a channel suppressed by \diag{W200} leaves
$\chanres{\st}{d}$ at its baseline value, and a bump from sources on another
channel is then suppressed by \cref{a:resolvable}. \Cref{ob:g8} states the
correction.
\end{remark}
